\begin{frame}[fragile]{Encodage basique}
Le XML repose sur trois principes~:
\begin{itemize}
    \item Les balises sont appelées \textbf{éléments};
    \item Les éléments sont imbriqués;
    \item L'imbrication est sous forme d'arborescence
\end{itemize}
\vspace{1em}

\begin{tabular}{c}
    \begin{lstlisting}[numbers=none]
    <elementParent>
        <elementEnfant> Texte </elementEnfant>
    </elementParent>
    \end{lstlisting}
\end{tabular}

\end{frame}

\begin{frame}[fragile]{Encodage basique}
\begin{center}
    \begin{tabular}{c}
    \begin{lstlisting}[numbers=none]
    <Paragraphe>
        <Phrase> 
            <mot> Toto </mot>
        </Phrase>
    </Paragraphe>
    \end{lstlisting}
\end{tabular}
\end{center}
    
\end{frame}

\begin{frame}{Éléments et attributs}
\begin{block}{}
    \textbf{Élément =} Le contenu. Doit toujours être entouré de deux balises du même nom.\\
    Ex: \textcolor{blue}{<paragraphe></paragraphe>} \textrightarrow{} l'élément est le paragraphe.\\

    Les éléments ne doivent pas se chevaucher.\\
    Ex: \textcolor{blue}{<paragraphe>}\textcolor{BurntOrange}{<mot>}\textcolor{blue}{</paragraphe>}\textcolor{BurntOrange}{</mot>}
\end{block}

\begin{block}{}
    \textbf{Attribut =} Caractéristique du contenu. Ne peut pas réutilisé deux fois le même attribut.\\ 
    Ex: \textcolor{blue}{<paragraphe} \textcolor{red}{numero="1"}\textcolor{blue}{></paragraphe>} \textrightarrow{} c'est le premier paragraphe.
\end{block}
    
\end{frame}

\begin{frame}{Une structure bien formée}

Un encodage XML est dit \textcolor{red}{\textbf{bien formé}} s'il respecte les règles de balisage ci-dessous~:

\begin{itemize}
    \item Chaque balise doit avoir une balise de fin (une balise ouvrante et une balise fermante)~;
    \item Les éléments ne doivent pas se chevaucher et doivent bien s'imbriquer~;
    \item Un élément ne peut avoir qu'un seul même attribut~;
    \item Un seul élément racine (on doit avoir un seul élément qui englobe tous les autres).
\end{itemize}

\end{frame}


\begin{frame}[fragile]{Déclaration XML}
\begin{lstlisting}
        <?xml version="1.0" encoding="UTF-8"?>
    \end{lstlisting}

\begin{itemize}
    \item La déclaration XML s'écrit en début de document
    \item Elle indique la version utilisée de XML (toujours la version 1.0) et l'encodage du texte.
\end{itemize}
\end{frame}


\begin{frame}[fragile]{Autres commandes à connaître}
\begin{block}{Commenter le texte}
    \begin{lstlisting}
        <!-- on écrit son commentaire -->
    \end{lstlisting}
\end{block}

\begin{block}{Les entités}

\begin{verbatim}
    &entité;
\end{verbatim}
Permet d'utiliser des caractères interdits en XML comme le \&.
    
    
\end{block}
    
\end{frame}

\begin{frame}{Exercices}
    \url{https://pad.libreon.fr/eoi6u4UjRhO-XGzRv_bPzg?both}
\end{frame}

